\documentclass{article}
\usepackage[UTF8]{ctex}
\usepackage{amsmath,amssymb,amsfonts}
\usepackage{enumitem}
\usepackage{geometry}
\geometry{a4paper, margin=1in}
% === 错题本排版优化（不参与题目标识） ===
\setlength{\parindent}{0pt}
\setlength{\parskip}{0.35em plus 0.10em minus 0.08em}
\setlength{\emergencystretch}{3em}
\setlist{itemsep=0.25em, topsep=0.40em, parsep=0pt}
\allowdisplaybreaks[2]
% === 错题本排版优化结束 ===
\begin{document}
\section*{中国人民大学2024年高等代数真题}
\section*{16．中国人民大学2024年高等代数试卷}
一．（25 分）填空题．\\

\subsection{题目1}
1．（5 分）\(\left|\begin{array}{ccccc}1 & 0 & 0 & 0 & 0 \\ 2 & 3 & 0 & 0 & 0 \\ 4 & 5 & 6 & 0 & 0 \\ 7 & 8 & 9 & 1 & -2 \\ 10 & 11 & 12 & 3 & 4\end{array}\right|=\)\\

\subsection{题目2}
2．（10 分）\(\left|\begin{array}{ccccc}0 & 1 & 0 & \cdots & 0 \\ 0 & 0 & 2 & \cdots & 0 \\ \vdots & \vdots & \vdots & & \vdots \\ 0 & 0 & 0 & \cdots & n-1 \\ n & 0 & 0 & \cdots & 0\end{array}\right|=\)\\

\subsection{题目3}
3．（10 分）\(\left|\begin{array}{cccc}1 & 1 & 1 & 1 \\ a & b & c & d \\ a^{2} & b^{2} & c^{2} & d^{2} \\ a^{4} & b^{4} & c^{4} & d^{4}\end{array}\right|=\)

二．（15 分）证明：设 \(A\) 是 \(n\) 阶矩阵，若有正整数 \(k\) ，使得 \(\operatorname{rank}\left(A^{k}\right)=\operatorname{rank}\left(A^{k+1}\right)\) ，则

\[
\operatorname{rank}\left(A^{k}\right)=\operatorname{rank}\left(A^{k+j}\right)(j=1,2,3, \cdots) .
\]

三．（15 分）证明：秩为 \(r\) 的矩阵 \(A\) 可分解为 \(A=C R\) ，其中 \(C, R\) 分别为列，行满秩矩阵，秩均为 \(r\) ，且若有两种这样的分解 \(A=C R=C_{1} R_{1}\) ，则存在可逆矩阵 \(P\) ，使 \(C P=C_{1}, P^{-1} R=R_{1}\) ．

四．（15 分）证明 Sylvester 不等式：设 \(A, C\) 为 \(m \times n, n \times q\) 矩阵，则有

\[
\operatorname{rank}(A)+\operatorname{rank}(C) \leq \operatorname{rank}(A C)+n
\]

五。（25分）设 \(V\) 是实数域 \(\mathbb{R}\) 上的向量空间，\(\alpha_{1}, \alpha_{2}, \cdots, \alpha_{n}\) 是其一组基，于是由

\[
\phi\left(\alpha_{i}\right)=\alpha_{i+1}(i=1, \cdots, n-1), \phi\left(\alpha_{n}\right)=0 .
\]

定义了一个 \(V\) 上的线性映射。\\

\subsection{题目4}
1．（5 分）写出 \(\phi\) 在 \(\alpha_{1}, \alpha_{2}, \cdots, \alpha_{n}\) 下的矩阵表示。\\

\subsection{题目5}
2．（5 分）证明：\(\phi^{n}=0, \phi^{n-1} \neq 0\) 。\\

\subsection{题目6}
3．（10 分）假设 \(\psi\) 是 \(V\) 上的线性映射，满足 \(\psi^{n}=0, \psi^{n-1} \neq 0\) ，则存在 \(V\) 的一组基使得 \(\psi\) 的矩阵表示与 \(\phi\) 在第1问中的矩阵表示相同。\\

\subsection{题目7}
4．（5 分）证明：若 \(n\) 阶实方阵 \(M, N\) 满足 \(M^{n}=N^{n}=O, M^{n-1} \neq O \neq N^{n-1}\) ，则 \(M\) 与 \(N\) 相似。

六。（30分）设向量空间 \(V\) 的线性变换 \(\phi\) 满足 \(\phi^{2}=\phi\) ，则称 \(\phi\) 为投影映射。证明：\\

\subsection{题目8}
1．（5 分）\(V\) 中向量 \(\beta\) 属于 \(\phi\) 的像集当且仅当 \(\phi(\beta)=\beta\) 。\\

\subsection{题目9}
2．（10 分）\(V=\phi(V) \oplus \operatorname{Ker}(\phi)\) ，且 \(V\) 中任何一个向量可以直和分解为 \(\alpha=\phi(\alpha)+(\alpha-\phi(\alpha))\) 。\\

\subsection{题目10}
3．（10 分）对任意直和分解 \(V=V_{1} \oplus V_{0}\) ，存在唯一的投影映射 \(\psi\) 使得 \(V_{1}=\psi(V), V_{0}=\operatorname{Ker}(\psi)\) ．\\

\subsection{题目11}
4．（5 分）每个投影映射都有矩阵表示 \(\left(\begin{array}{cc}I & O \\ O & O\end{array}\right)\) 。

七．（25 分）设 \(A\) 是 \(n\) 阶复方阵，计算所有与 \(A\) 可交换的复方阵，并写出计算过程．
\end{document}
